\subsection{DUAL OF A QUOTIENT MODULE. DUAL OF A DIRECT SUM. DUAL BASES}

We apply Theorem 1 of no. 1 to the case where \(\mathbf{F} = \mathbf{A}_{\mathrm{s}}\) :

PROPOSITION 9. Let \(\mathbf{E}'\) , \(\mathbf{E}\) , \(\mathbf{E}''\) be A-modules and

\[
\mathrm{E} ^ {\prime} \stackrel {u} {\longrightarrow} \mathrm{E} \stackrel {v} {\longrightarrow} \mathrm{E} ^ {\prime \prime} \longrightarrow 0\tag{21}
\]

an exact sequence of linear mappings. Then the sequence of transpose mappings

\[
0 \longrightarrow \mathrm{E} ^ {\prime \prime *} \stackrel {t _ {v}} {\longrightarrow} \mathrm{E} ^ {*} \stackrel {t _ {u}} {\longrightarrow} \mathrm{E} ^ {\prime *} \tag {1}
\]

is exact.

COROLLARY. Let M be a submodule of an A-module E and \(\phi: \mathbf{E} \to \mathbf{E} / \mathbf{M}\) the canonical homomorphism. Then \(^t\phi\) is an isomorphism of the dual of E/M onto the submodule M' of E* orthogonal to M.

If \(j: \mathbf{M} \to \mathbf{E}\) is the canonical injection, the kernel of \(j\) is by definition the orthogonal of \(\mathbf{M}\) in \(\mathbf{E}^*\) .

Moreover, in the notation of the corollary, an injective homomorphism \(E^{*}/M^{\prime} \rightarrow M^{*}\) is obtained from \(^{t}j\) when passing to the quotient.

PROPOSITION 10. Let \((\mathbf{E}_t)_{t \in I}\) be a family of A-modules and for all \(t \in I\) let \(j_t: \mathbf{E}_t \to \mathbf{E} = \bigoplus_{t \in I} \mathbf{E}_t\) be the canonical injection. Then the product mapping \(x^* \mapsto (t j_t(x^*))\) is an isomorphism of the dual \(\mathbf{E}^*\) of \(\mathbf{E}\) onto the product \(\prod_{t \in I} \mathbf{E}_t^*\) .

This is a particular case of §1, no. 6, Corollary 1 to Proposition 6, applied to the case where \(\prod_{\lambda} F_{\lambda} = A_{s}\) .

If, by means of the canonical injections \(j_{i}\) , the \(E_{i}\) are identified with submodules of their direct sum \(E\) and if, by means of the product mapping \(x^{*} \mapsto (t j_{i}(x^{*}))\) , \(E^{*}\) is identified with \(\prod_{i \in I} E_{i}^{*}\) , it can then be said that \(\prod_{i \in I} E_{i}^{*}\) is the dual of \(\bigoplus_{i \in I} E_{i}\) , the canonical bilinear form being given by

\[
\langle (x _ {i}), (x _ {i} ^ {*}) \rangle = \sum_ {i \in 1} \langle x _ {i}, x _ {i} ^ {*} \rangle .\tag{22}
\]

COROLLARY. Let M, N be two supplementary submodules in an A-module E and \(p:E\to M\) , \(q:E\to N\) the corresponding projectors; then \({}^{tp}+{}^{tq}:M^{*}\oplus N^{*}\to E^{*}\) is an isomorphism and \({}^{tp}\) (resp. \({}^{tq}\) ) is an isomorphism of \(M^{*}\) (resp. \(N^{*}\) ) onto the submodule of \(E^{*}\) orthogonal to N (resp. M). Moreover, if \(i:M\to E\) and \(j:N\to E\) are the canonical injections, \({}^{tp}\circ{}^{ti}\) and \({}^{tq}\circ{}^{tj}\) are the projectors \(E^{*}\to{}^{tp}(M^{*})\) , \(E^{*}\to{}^{tq}(N^{*})\) corresponding to the decomposition of \(E^{*}\) as the direct sum of \({}^{tp}(M^{*})\) and \({}^{tq}(N^{*})\) .

\(p \circ i = 1_{\mathbf{M}}, \quad q \circ j = 1_{\mathbf{N}}, \quad p \circ j = q \circ i = 0, \quad i \circ p + j \circ q = 1_{\mathbf{E}},\) whence, by transposition (no. 5, formulae (16), (17) and (18)), \(^{ti} \circ ^{tp} = 1_{\mathbf{M}^*}, \; ^{tj} \circ ^{tq} = 1_{\mathbf{N}^*}, \; ^{tj} \circ ^{tp} = ^{ti} \circ ^{tq} = 0. \; ^{tp} \circ ^{ti} + ^{tq} \circ ^{tj} = 1_{\mathbf{E}^*}\) and the proposition follows from § 1, no. 6, Corollary 2 to Proposition 6.

Under the hypotheses of the Corollary, \(M^{*}\) (resp. \(N^{*}\) ) is often identified with the orthogonal \(^{tp}(M^{*})\) (resp. \(^{t}q(N^{*})\) ) of N (resp. M) in \(E^{*}\) , thus identifying every linear form u on M (resp. N) with the linear form on E extending u and which is zero on N (resp. M).

When an A-module E admits a basis \((e_{t})_{t\in\mathbb{T}}\) , it has been seen that giving this basis defines canonically an isomorphism \(u:A_{s}^{\mathrm{(T)}}\to E\) . By virtue of Proposition 10 and no. 3, Proposition 5, the dual of \(\mathbf{A}_{s}^{\mathrm{(T)}}\) is canonically identified with the product \(A_{d}^{T}\) ; consider the contragredient isomorphism \(\check{u}:A_{d}^{T}\to E^{*}\) . If, for all \(t\in T,f_{t}\) is the element of \(A_{d}^{T}\) all of whose projections are zero with the exception of that of index t, which is equal to 1, and if we write \(e_{t}^{*}=\check{u}(f_{t})\) , the elements \(e_{t}^{*}\) of \(E^{*}\) are, by (19) and (22), characterized by the relations

\[
\langle e _ {t}, e _ {t ^ {\prime}} ^ {*} \rangle = \left\{ \begin{array}{l l} 0 & \text { for } t ^ {\prime} \neq t \\ 1 & \text { for } t ^ {\prime} = t. \end{array} \right.\tag{23}
\]

It amounts to the same to say that, for all \(x = \sum_{t \in T} \xi_t e_t \in E\) , \(e_t^*(x) = \xi_t\) ; also \(e_t^*\) is called the coordinate form of index \(t\) on \(E\) . It follows from (23) that \((e_t^*)\) is a free system in \(E^*\) .

In particular, if T is finite, the \(e_{t}^{*}\) form a basis of \(E^{*}\) , the \(f_{t}\) then forming the canonical basis of \(A_{d}^{T}\) . Hence:

PROPOSITION 11. The dual of a free module with a basis of \(n\) elements is a free module with a basis of \(n\) elements.

Note that the dual of a free module with an infinite basis is not necessarily a free module (VII, § 3, Exercise 10).

DEFINITION 7. If E is a free module with a finite basis \((e_{t})\) , the basis \((e_{t}^{*})\) of the dual \(E^{*}\) of E defined by relations (23) is called the dual basis of \((e_{t})\) .

Relations (23) can also be written in the form

\[
\langle e _ {t}, e _ {t ^ {\prime}} ^ {*} \rangle = \delta_ {t t ^ {\prime}}\tag{24}
\]

where \(\delta_{tt'}\) is the Kronecker symbol on \(\mathbf{T} \times \mathbf{T}\) .

Note that if T is finite and \((e_{t}^{*})\) is the dual basis of \((e_{t})\) , then, for \(x = \sum_{t \in T} \xi_{t} e_{t} \in E\) , \(x^{*} = \sum_{t \in T} \xi_{t}^{*} e_{t}^{*} \in E^{*}\) ,

\[
\langle x, x ^ {*} \rangle = \sum_ {t \in T} \xi_ {t} \xi_ {t} ^ {*}.\tag{25}
\]

The dual basis of a finite basis of a right A-module is of course defined similarly.

COROLLARY. The dual of a finitely generated projective module is a finitely generated projective module.

A finitely generated projective left A-module can be identified with a direct factor M of a free A-module \(A_{s}^{n}\) with a finite basis (no. 2, Corollary 1 to Proposition 4). Then (Proposition 11 and Corollary to Proposition 10) \(M^{*}\) is isomorphic to a direct factor of \(A_{d}^{n}\) , whence the corollary.

PROPOSITION 12. Let E be an A-module and \((a_{t})_{t\in \mathbf{T}}\) a generating system of E. The following conditions are equivalent:

\begin{enumerate}
\def\labelenumi{(\alph{enumi})}
\item
  E is a projective A-module.
\item
  There exists a family \((a_{t}^{*})_{t\in \mathbf{T}}\) of linear forms on \(\mathbf{E}\) such that, for all \(x\in \mathbf{E}\) , the family \((\langle x,a_t^*\rangle)_{t\in \mathbf{T}}\) has finite support and
\end{enumerate}

\[
x = \sum_ {t \in \mathbf {T}} \langle x, a _ {t} ^ {*} \rangle a _ {t}.\tag{26}
\]

There exists a surjective homomorphism \(u: \mathbf{L} \to \mathbf{E}\) , where \(\mathbf{L} = \mathbf{A}_s^{(\mathrm{T})}\) , such that if \((e_t)_{t \in \mathbb{T}}\) is the canonical basis of \(\mathbf{L}\) then \(u(e_t) = a_t\) (§ 1, no. 11, Proposition 17); for \(\mathbf{E}\) to be projective, it is necessary and sufficient that there exist a linear mapping \(v: \mathbf{E} \to \mathbf{L}\) such that \(u \circ v = 1_{\mathbb{E}}\) (no. 2, Proposition 4 and § 1, no. 9, Proposition 15). If such a mapping exists and we write \(^tv(e_t^*) = a_t^*\) , then \(\langle x, a_t^* \rangle = \langle x, ^tv(e_t^*) \rangle = \langle v(x), e_t^* \rangle\) , hence the family \((\langle x, a_t^* \rangle)\) has finite support and \(x = u\left(\sum_{t \in \mathbf{T}} \langle (x), e_t^* \rangle e_t\right) = \sum_{t \in \mathbf{T}} \langle x, a_t^* \rangle a_t\) for all \(x \in \mathbf{E}\) . Conversely, if condition (b) of the statement is fulfilled, the sum \(\sum_{t \in \mathbf{T}} \langle x, a_t^* \rangle e_t\) is defined for all \(x \in \mathbf{E}\) and \(x \to \sum_{t \in \mathbf{T}} \langle x, a_t^* \rangle e_t\) is a linear mapping \(v: \mathbf{E} \to \mathbf{L}\) such that \(u \circ v = 1_{\mathbf{E}}\) .

